%\section{Ottica Geometrica: Rifrazione}
\begin{frame}{La Rifrazione}
    \begin{itemize}
        \item Quando la luce passa tra due mezzi con $n$ diverso, cambia velocità e devia la traiettoria
        \item Il raggio si avvicina alla normale entrando in un mezzo più denso ($n_2 > n_1$)
        \item Il raggio si allontana dalla normale entrando in un mezzo meno denso ($n_2 < n_1$)
        \item Esempio immediato: un oggetto immerso in acqua appare spezzato o spostato
    \end{itemize}
\end{frame}

%========================================================

\begin{frame}{La Legge di Snell-Descartes}
    $$n_1 \sin\theta_1 = n_2 \sin\theta_2$$
    \begin{itemize}
        \item $n_1$, $n_2$: indici di rifrazione dei due mezzi
        \item $\theta_1$: angolo di incidenza (rispetto alla normale)
        \item $\theta_2$: angolo di rifrazione (rispetto alla normale)
        \item Il prodotto $n \sin\theta$ si \textbf{conserva} attraversando l'interfaccia
        \item Raggio incidente, normale e raggio rifratto: \textbf{stesso piano}
    \end{itemize}
\end{frame}


%========================================================

\begin{frame}{Riflessione Totale Interna}
    \begin{itemize}
        \item  Si verifica quando la luce passa da un mezzo più denso a uno meno denso ($n_1 > n_2$)
        \item Esiste un angolo critico $\theta_c$ oltre il quale la rifrazione è impossibile:
        $$\sin\theta_c = \frac{n_2}{n_1}$$
        \item Per $\theta_i > \theta_c$: tutta la luce viene riflessa, nessuna trasmessa
        \item L'interfaccia diventa uno specchio perfetto
    \end{itemize}
\end{frame}


%========================================================



